%% =====================================================================
%%  B1-T-20-02 -- what B1 contributes to "The Dirty Way"
%%  -------------------------------------------------------------------
%%  LAYER A: APPEND-ONLY. Written once at the entry's write, then left
%%  alone. If the source has more to say later, add a bullet -- do not
%%  rewrite. Material found ONLY here goes in the `unique` slot with
%%  @unique set to yes, which makes it undeletable (rule R10).
%% =====================================================================
%% @kind: source
%% @id: B1-T-20-02
%% @book: B1
%% @topic: T-20-02
%% @locator: Tactic 13 (ch. XIV), pp. 161--166
%% @status: distilled
%% @unique: yes
%% @integrated: yes
%% @updated: 2026-09-19

\dossier{B1}{T-20-02}{\bookshort{B1}}{Tactic 13 (ch. XIV), pp. 161--166}

\begin{distilled}
  \item When the front door closes, the source admits the indirect route: reach the target through spouse, friends, household, loyalties.
  \item The worked example is a salesman with a ruined name who made the customer's wife the seller, so the purchase arrived as the household's own idea.
  \item The lever is borrowed trust: refusal aimed at the operator costs nothing; refusal aimed at a loved one costs intimacy.
  \item The source itself files this as manipulation against the will --- the con artist's road --- and its closing chapters decline to extend the course.
  \item Much of the conduct described is criminal in most jurisdictions; the tactic is admitted for recognition, not practice.
\end{distilled}
\begin{unique}
  \item The classified-admissions framing: the chapter is written as the book's own warning label for the road it refuses to teach further.
\end{unique}
